\section{把很多人的偏好合起来}
\label{sec:17-Constrained-Guarantees/05-Learning-from-Comparisons/03}

质检抽了三百条比较发给三个标注员重做，一致率 \(68\%\)。项目组的处理办法是加人：以后每条比较发给三个人，取多数。

这个办法在单条比较上确实降低了噪声。问题出现在把结果拼起来的时候。抽检里有一组三个回答 \(A,B,C\)，三个标注员各自给出了一个完整的、内部自洽的偏好序：甲说 \(A\succ B\succ C\)，乙说 \(B\succ C\succ A\)，丙说 \(C\succ A\succ B\)。

按多数：比 \(A\) 与 \(B\)，甲丙都选 \(A\)，二比一，\(A\) 胜；比 \(B\) 与 \(C\)，甲乙都选 \(B\)，\(B\) 胜；比 \(C\) 与 \(A\)，乙丙都选 \(C\)，\(C\) 胜。于是 \(A\succ B\succ C\succ A\)。

三个人都没有自相矛盾，三次多数表决都没有算错，合起来的结果却排不出序。这不是标注质量问题，加到九个人、九十个人都可能出现同样的形状。

\begin{center}
\begin{tikzpicture}[x=1cm,y=1cm,font=\small,>=stealth,
  itm/.style={draw,circle,minimum size=0.6cm,inner sep=0pt,fill=blue!8}]
  % 三位标注员的偏好序
  \foreach \i/\p/\q/\r in {0/A/B/C, 1/B/C/A, 2/C/A/B} {
    \begin{scope}[shift={(\i*1.3,0)}]
      \draw[black!50,fill=blue!7] (0,1.8) rectangle (1.0,2.4);
      \node at (0.5,2.1) {\(\p\)};
      \draw[black!50,fill=blue!7] (0,1.1) rectangle (1.0,1.7);
      \node at (0.5,1.4) {\(\q\)};
      \draw[black!50,fill=blue!7] (0,0.4) rectangle (1.0,1.0);
      \node at (0.5,0.7) {\(\r\)};
    \end{scope}
  }
  \node[font=\footnotesize] at (0.5,2.7) {甲};
  \node[font=\footnotesize] at (1.8,2.7) {乙};
  \node[font=\footnotesize] at (3.1,2.7) {丙};
  \node[font=\footnotesize,anchor=north] at (1.8,0.15) {三个内部自洽的偏好序};
  \node[font=\footnotesize,anchor=east] at (-0.1,2.1) {更好};
  \node[font=\footnotesize,anchor=east] at (-0.1,0.7) {更差};
  % 循环
  \begin{scope}[xshift=6.0cm,yshift=0.35cm]
    \node[itm] (a) at (0,0.2) {\(A\)};
    \node[itm] (b) at (2.5,0.2) {\(B\)};
    \node[itm] (c) at (1.25,2.0) {\(C\)};
    \draw[->,black!70,thick] (a) -- (b);
    \draw[->,black!70,thick] (b) -- (c);
    \draw[->,black!70,thick] (c) -- (a);
    \node[font=\footnotesize,anchor=north] at (1.25,0.0) {\(2:1\)};
    \node[font=\footnotesize,anchor=west] at (2.05,1.25) {\(2:1\)};
    \node[font=\footnotesize,anchor=east] at (0.45,1.25) {\(2:1\)};
    \node[font=\footnotesize,anchor=north] at (1.25,-0.5) {两两多数给出一个环};
  \end{scope}
\end{tikzpicture}

\emph{左边三列每一列都是一个完整的序，右边那个环是它们两两取多数的结果。个体传递不蕴涵集体传递，所以``按多数意见排序''这个操作本身可能没有定义。}
\end{center}

这个环叫 Condorcet 悖论。它与上一节那个非传递循环形状一样，来源却不同：那里的循环来自单个标注员在不同维度上换了注意力，这里的循环来自不同人之间的分歧被逐对压缩。同一个环，\hyperref[sec:04-Discrete-Mathematics/03-Graph-Theory/02]{``路径、环与连通性''}里的判断在这里成立——有向环的存在说明这张图不能被拓扑排序，而排序正是我们要的东西。

\subsection{不是多数规则的问题}

自然的反应是换一个更聪明的规则。Borda 计分、逐轮淘汰、加权投票，都有人提过。Arrow 的定理说明这条路不通：不是当前规则不够好，是符合几条基本要求的规则不存在。

\begin{theorem}[Arrow]
设选项数不少于三，个体数有限。若聚合规则对每一组个体偏好序都输出一个完整且传递的社会偏好序，且满足帕累托一致与无关方案独立性，则该规则必然是独裁的——存在某个个体，社会序永远等于他的序。
\end{theorem}

证明这里不给，它是一个巧妙但不短的组合论证。真正要交代清楚的是三条前提各挡住什么失效模式，因为定理的用处全在这里。

\textbf{帕累托一致}要求：若每个人都认为 \(A\) 优于 \(B\)，社会序也必须如此。它挡住的是规则与输入脱钩——一个恒定输出某个固定序的规则满足其余所有条件，但它根本没在聚合。这条要求把``规则必须真的听取偏好''写成了形式条件。

\textbf{无关方案独立性}要求：社会对 \(A\) 与 \(B\) 的排序，只能依赖各人对 \(A\) 与 \(B\) 的相对排序，与他们怎么看 \(C\) 无关。它挡住的是议程操纵——若这条不成立，往候选集里塞一个明显更差的选项就能翻转 \(A\) 与 \(B\) 的胜负。Borda 计分正是放弃了这一条：它给每个选项按名次赋分，而加入一个新选项会改变所有名次。放弃 IIA 换来的是可用的规则，让出的是对无关选项的免疫。

\textbf{无独裁者}挡住的是最平凡的解法。把某一个人的序直接当作社会序，帕累托一致成立，无关方案独立性也成立——它满足其余全部要求，所以必须被单独排除，否则定理无内容。

还有一条隐含前提值得点出来：\textbf{定义域无限制}，即规则必须对任意一组个体偏好都有定义。这是唯一一个有实际逃生口的前提。若能确定所有人的偏好是单峰的——存在一个共同的一维坐标，每个人有一个最偏好的点，离它越远越差——那么多数规则就不会产生环，中位数投票者的最优点就是稳定的胜者。为此要付出的是一句论证：单峰性得成立，而这是一句关于人群的实质断言，不能靠假设省事。同样地，若把``输出一个完整传递序''放松成``只选出一个胜者''，Arrow 的结论不直接适用，但另一组关于策略操纵的定理会接上来。前提松掉一条，麻烦换一个地方出现，没有一处是白让的。

\subsection{损失函数已经替你选了一个聚合规则}

回到奖励模型。上一节的训练目标是把所有比较的负对数似然加起来：
\[
\sum_{k}-\log\sigma\bigl(r_\theta(x_k,y_{w,k})-r_\theta(x_k,y_{l,k})\bigr),
\]
求和跑遍全部标注记录，不区分是谁标的。这个求和看上去只是记账，它其实完整地指定了一条聚合规则。

规则的性质可以逐条对照。它\textbf{必然输出传递序}，因为 \(r_\theta\) 取值在实数上，实数的序是全序——所以遇到上面那个环时，模型不会拒绝作答，它会给出一组分数，让三个胜率都贴近 \(1/2\)。环被抹平成三个平局，而分歧的存在从输出里消失了。它\textbf{不满足无关方案独立性}：\(A\) 与 \(B\) 的估计分差由整张比较图上的全部数据共同决定，往图里加入一个 \(C\) 并采集若干含 \(C\) 的比较，会改变 \(A\) 与 \(B\) 的分差。放弃 IIA 正是它能给出传递输出而不沦为独裁的原因，这与 Borda 走的是同一条路。

由此可以读出一个不太舒服的推论：\textbf{哪一对被采样得多，哪一对在聚合结果里的权重就大}。采样设计因此是聚合规则的一部分，而它通常由工程便利决定——同一个 prompt 下比两个回答、热门 prompt 采得多、难以判断的比较对被标注员跳过。这些决定没有一条是在讨论``该怎么聚合偏好''时做出的，但它们的效果正是在做这件事。

Arrow 的结论在这里不是提供一个禁令，而是取消一个幻想：不存在一个在所有情形下都合理的聚合规则，所以``众包标注取多数''不是中立操作，也不存在一个更中立的替代品。能做的是把已经做出的选择写下来——用了哪条规则、它放弃了哪条要求、放弃之后会以什么形式出错。这正是\hyperref[sec:18-Synthesis/01-System-Lifecycle/01]{``把诉求写成目标函数与约束''}里那步翻译，只不过这次被翻译的不是一个人的诉求，是一群人不相容的诉求。而\hyperref[sec:11-Mathematical-Statistics/07-Statistical-Decision-and-Reliable-Prediction/01]{``状态、行动、损失与风险构成一个决策问题''}里那套记账方式提示了一个实际做法：与其把分歧压平，不如把标注员身份当作一个变量留在模型里，让分歧显式地进入报告——这不解决 Arrow 的问题，但至少不假装它不存在。

\subsection{这个单元留下什么}

三节沿同一条线走完。人对``哪个更好''的判断比对``值几分''的判断稳定，这是经验规律，也是从比较学习的全部理由。而比较能定出的东西有明确边界：只有差值，没有水平；只在比较图的连通分支之内，跨分支的关系无从谈起；把差值压进一个实数分数，已经假设了偏好可传递。

把分数换成一个以输入为条件的函数，这三条限制一条不少，只是更难察觉。它们的后果是具体的：奖励的绝对值是未定义的量，任何依赖它的下游做法都在使用一个换个随机种子就会变的数；而优化过程会把策略推出奖励模型有证据的区域，于是奖励继续涨、被奖励衡量的东西开始跌——这一条是被反复观察到的经验规律，分岔点在哪里没有可计算的答案。

而把多人的比较不加区分地加进同一个损失，等于选定了一条偏好聚合规则。Arrow 的定理说明这个选择躲不掉：满足帕累托一致、无关方案独立、无独裁的传递聚合规则不存在，所以任何可用的规则都已经在某处让步。奖励模型让的是无关方案独立性，让步的形式是采样设计悄悄变成了权重设计。

这条线索与本部分其余单元共享同一个形状：几件都想要的好事被一条恒等式、一条界或一条定理绑在一起，于是必须指定要哪一件。这里被绑住的是``尊重每个人的判断''、``给出一个可排序的结果''、以及``结果不受无关选项影响''——三者不能兼得，而这与数据量、模型规模、标注预算都无关。
